\summarychapter{24}{A Wild Boar in the Vineyard}{Martin Luther and Protestantism}

\begin{inbrief}
The differences between Catholics and Protestants go back to the age of Luther, the Reformation
(1517--1648). Protestantism gave new answers to four old Catholic questions, and those answers first
took shape in one German monk's struggle of soul and his conflict with Rome. The chapter follows
Luther from the monastery to the Ninety-five Theses, the break with the pope, the Diet of Worms and
the making of a Lutheran church.
\question{What is Protestantism? How is a person saved, where does religious authority lie, what is
the church, and what is the essence of Christian living?}
\end{inbrief}

In the summer of 1520 a papal bull went through Germany in search of Martin Luther: ``Arise, O Lord,
and judge thy cause. A wild boar has invaded thy vineyard.'' It condemned forty-one of his beliefs and
gave him sixty days to repent. Luther received it on October 10. When the time ran out in early
December he led his students outside Wittenberg and burned the canon law, some medieval theologians
and, as if by afterthought, the bull itself: ``They have burned my books; I burn theirs.''

\ssection{The Meaning of Protestantism}
Ernst Troeltsch called Protestantism a ``modification of Catholicism'' in which the Catholic problems
remain but receive different solutions. The four questions it answered anew concern salvation,
authority, the church and Christian living. Reformers all over Europe reached similar convictions,
but the answers emerged first in Luther.

Born in 1483, son of a Saxon miner, Luther meant to be a lawyer until in 1505 a bolt of lightning
near Stotternheim threw him to the ground and he cried, ``St.~Anne, save me, and I'll become a
monk.'' Two weeks later he entered the Augustinian monastery at Erfurt. No monk was more zealous: he
fasted for days, slept without blankets in winter, and said that if ever a monk got to heaven by
monkery it would have been he. He was driven by a sense of his own sinfulness and of God's majesty.
Saying his first mass he was terror-stricken, and no penance or advice quieted his conviction that he
was doomed. ``I do not love God! I hate him!'' he once burst out.

He found what he sought in Scripture, as professor of biblical studies at the new university of
Wittenberg. Christ's cry of forsakenness on the cross made him ask whether Christ had shared
humanity's estrangement in order to bear the punishment of sin. In 1515, pondering Romans 1:17, ``The
just shall live by faith,'' he grasped that the justice of God is the righteousness by which God in
sheer mercy justifies us through faith, and he felt reborn, as if he had gone ``through open doors
into paradise.'' A person is saved only by faith in the merit of Christ's sacrifice. This is
justification by faith alone, and it clashed with Rome's justification by faith and good works. The consequences reached far. If faith in Christ is all that saves, a faith fed by the Word
written and preached, then the intercession of priests, monks, masses and prayers to the saints are
unnecessary, and the mediation of the Roman church crumbles.

\ssection{Luther's Attack on Papal Authority}
A flagrant abuse of church finance drove him further. Indulgences, begun in the crusades, remained a
favorite source of papal income. In return for a meritorious work, often a gift of money, the church
released the sinner from acts of penance by drawing on the ``treasury of merits.'' Preachers made
them seem a kind of magic and ignored sorrow for sin. In 1517 the Dominican John Tetzel was selling
them through Germany for the building of St.~Peter's in Rome, promising release even for souls in
purgatory: ``As soon as the coin in the coffer rings, the soul from purgatory springs.'' On October
31, 1517, Luther posted ninety-five theses for debate on the door of the Castle Church. The theses held that an indulgence takes away no guilt, has no effect in purgatory and lulls the buyer into false security.
That was the spark.

The German Dominicans denounced him to Rome, and a Vatican theologian declared criticism of
indulgences to be heresy. Luther began to demand proof from Scripture. In an eighteen-day debate at
Leipzig in 1519 John Eck drove him to say that a council may err and that articles of faith cannot be laid down by church or pope but only drawn from Scripture. This was his second conviction: the
Scriptures, not popes or councils, are the standard. Eck, noting the likeness to Hus, pressed Rome to
condemn him, and Luther appealed to the German people in three pamphlets of 1520. The \emph{Address to
the Nobility of the German Nation} called on the princes to correct abuses, strip bishops and abbots
of wealth and worldly power and create in effect a national German church. \emph{The Babylonian
Captivity of the Church} charged that Rome's sacramental system holds Christians captive and denies
them direct access to God by faith. A sacrament must be instituted by Christ and be exclusively
Christian; by that test only baptism and the Lord's Supper remain, and these belong to the community
of believers, not to an exclusive priesthood.

\begin{sidebox}{Important Theological Themes for Luther}
\emph{Justification by faith.} Catholics followed Augustine in teaching that righteousness is
imparted, an inner possession acquired over a lifetime. Luther taught imputation, a bookkeeping term:
God credits the righteousness of Jesus to the believer, who remains a sinner in experience and can
begin to become righteous only after being declared so. Protestants call the declaration justification and the growth that follows it sanctification. A table in the book, ``Summary of Righteousness,'' compares the two models:
\begin{sumtable}{L{0.47\linewidth}L{0.47\linewidth}}
\sumhead{Augustine (Catholic model)} & \sumhead{Luther (Protestant model)} \\ \midrule
The faithful are made righteous in the course of the spiritual pilgrimage &
  Believers are declared righteous; only then can they grow holier in experience \\
Righteousness comes gradually, as a process & Righteousness comes at once, as an event \\
It is imparted & It is imputed \\
It is an inner state, infused or implanted & It is external to the believer \\
Being justified and being made holy are one thing & Being justified and being made holy are kept apart \\
\end{sumtable}
\emph{Law and Gospel.} Luther found two orientations throughout the Bible, and both promise blessing.
The Law's promise depends on keeping the covenant's conditions; it torments those who cannot, yet it
exposes spiritual bankruptcy and the need of grace. The Gospel's promise is greater and
unconditional, because God in mercy meets every condition. Gospel passages occur in both Testaments,
most of all in Isaiah, John and Paul, and not at all in James.

\emph{Theology of the cross.} To an eyewitness the crucifixion looked like God's defeat; only later
is it seen as his victory over evil. The cross reverses values: grace is for humble sinners who know
their incapacity, not for the disciplined and decent. God makes sinners saints and the dead alive.
Catholic theology, confident that man can ascend toward God, Luther called a theology of glory.
\end{sidebox}

Luther thus set aside the church as a sacred hierarchy under the pope and returned to the community
of believers in which all are priests. His third pamphlet, \emph{The Freedom of a Christian Man}, the
best introduction to his thought, argues that the inner freedom given by the certainty of faith leads
to good works: ``Good works do not make a man good, but a good man does good works.'' The essence of
Christian living is to serve God in one's calling, secular or ecclesiastical; all useful callings are
equally sacred. Monasticism was no longer necessary.

\ssection{Heretic, Outlaw, and Hero}
Leo X issued his bull in June 1520, and in January 1521 he declared Luther a heretic and expelled him
from the church. The young emperor Charles V, sworn to defend the church, summoned him to the Diet of
Worms. There Luther refused to recant: ``My conscience is captive to the Word of God. \dots{} Here I
stand, I cannot do otherwise.'' Charles declared him an outlaw. He was saved by Frederick the Wise of
Saxony, who hid him in the Wartburg Castle, where, disguised as ``Junker George,'' he spent nearly a
year and translated the New Testament into German.

Meanwhile the revolt spread. Towns removed statues and abandoned the mass, more radical reformers
appeared, and princes defied the ban. In 1522 Luther returned to Wittenberg and carried through a
reform that much of Germany copied. The office of bishop disappeared; clergy, monks and
nuns married, and he himself married a former nun, Katherine von Bora, so that a new figure appeared,
the married pastor with his family. He translated and revised the liturgy, gave the laity both bread
and wine as the Hussites had asked, and moved the center of worship from the sacrifice of the mass to
the preaching of the Word.

In 1524 the cost of his alliance with the princes appeared. The peasants, applying Christian freedom
to their social condition, rose against their lords with twelve demands, among them the abolition of
serfdom. Luther at first admitted the justice of their complaints, but when they turned violent he
urged the princes, in \emph{Against the Thievish and Murderous Hordes of Peasants}, to ``knock down,
strangle, and stab.'' In 1525 the revolt was crushed at a cost of perhaps one hundred thousand lives.
The survivors thought Luther a false prophet, and many returned to Catholicism or joined the
radicals. He held that the equality of all before God is spiritual, not social. That view pleased the
princes, many of whom became Lutheran partly because it let them control the church in their lands.
(A chart in the chapter shows the ``Major Traditions of the Reformation.'')

\ssection{Luther's Lasting Influence}
By 1530 leadership was passing from Luther. Still an outlaw, he could not attend the meeting at
Augsburg, where Philip Melanchthon, a young professor of Greek, wrote the Augsburg Confession, to which the Lutheran princes and theologians put their names; the emperor remained unmoved. Luther's last years are
hard to defend. He approved the bigamy of Philip of Hesse, abused reformers who disagreed with him as
he had once abused the papacy, and wrote of the Jews in terms that would be excessive in a Hitler. At
his death in 1546 he was, says Roland Bainton, an irascible and coarse old man. These faults do not
lessen an achievement that changed Christianity and Western civilization.

Charles V set out to crush the heresy. The Lutheran princes formed the Schmalkald League in 1531,
and from 1546 to 1555 there was intermittent civil war. The Peace of Augsburg (1555) let each prince
decide the religion of his subjects, permitted no Protestantism except Lutheranism, and required
Catholic bishops who turned Lutheran to give up their property. Lutheranism became a state religion
in large parts of Germany and spread to Scandinavia, and the individual had to believe what his
prince believed. Luther's greatest contribution was religious. A person is saved not by works but by
faith alone; authority lies not in the visible Roman institution but in the Word of God in the Bible;
the church is the whole community of believers, all of them priests; and Christian living is serving
God in any useful calling, ordained or lay. Any classical account of Protestantism echoes these four
answers.

\begin{keyfacts}
\keypeople{Martin Luther; John Tetzel; John Eck; Pope Leo X; Charles V; Frederick the Wise; Katherine
von Bora; Philip Melanchthon; Philip of Hesse.}
\keydates{%
  \kd{1483}{Birth of Martin Luther}
  \kd{1505}{Luther enters the Augustinian monastery at Erfurt}
  \kd{1515}{Luther grasps justification by faith from Romans 1:17}
  \kd{1517}{Luther posts the Ninety-five Theses (October 31)}
  \kd{1519}{Luther debates John Eck at Leipzig}
  \kd{1520}{Papal bull against Luther; his three great pamphlets; he burns the bull}
  \kd{1521}{Luther excommunicated; ``Here I stand'' at the Diet of Worms; Wartburg}
  \kd{1525}{The Peasants' Revolt crushed}
  \kd{1530}{Augsburg Confession drafted by Melanchthon}
  \kd{1546}{Death of Luther}
  \kd{1555}{Peace of Augsburg: each ruler decides the religion of his territory}}
\keyterms{Reformation; papal bull; justification by faith alone; imputed righteousness;
sanctification; Law and Gospel; theology of the cross; indulgences; authority of Scripture; priesthood
of all believers; calling; Augsburg Confession; Schmalkald League.}
\end{keyfacts}
